\begin{abcliste}
\abc Where the pulses overlap, the displacements add. The two pulses have the same shape but opposite signs, so the spring is straight when they cover each other completely, i.e.\ when their left ends (at $x_1 = \xaO$ and $x_2 = \xbO$) meet. The pulses approach each other at $2v$:
\begin{align}
 t &= \tF = \frac{\xb-\xa}{2\times\v} = \t \approx \result{\tP{0}{3}}
\end{align}
(Their fronts meet already at \SI{1.0}{\s}; then they only start to overlap.)

\abc At that moment the spring is straight, but its coils are moving fast: all the energy is $\result{\text{kinetic energy of the coils}}$. The energy is not lost: a moment later the pulses come out again and move on.
\end{abcliste}
